\begin{textsample}{POS Dim 2 – 2009\_c6 – Score -1.00 – 2009\_c6/enade\_2009\_biblio\_0e14e9ac55.txt}  \label{ex:f2_pos_303}
Capítulo 3 Análise Técnica da Prova Este capítulo tem por objetivo apresentar o desempenho dos estudantes de BIBLIOTECONOMIA no ENADE/2009.
Para isso, foram calculadas as estatísticas básicas da prova como um todo, bem como as estatísticas das questões relacionadas à Formação Geral e ao Componente Específico.
Nas Tabelas, são apresentadas as seguintes estatísticas : total da população, de presentes, média, erro-padrão da média, desvio-padrão, mediana, nota máxima, nota mínima e coeficiente de assimetria ( ver definição em 1.4.3 ).
Nas estatísticas das questões discursivas, foram detalhadas também as notas nulas : em branco, por protesto ou por insuficiência.
As estatísticas apresentadas neste capítulo contemplam, separadamente, os ingressantes, os concluintes e o total de estudantes, e foram calculadas tendo -se em vista as seguintes agregações : ( a ) Região Geográfica e País e ( b ) a Categoria Administrativa e a Organização Acadêmica.
As notas brutas das questões discursivas em cada uma das partes, Formação Geral/Componente Específico, foram calculadas a partir da média simples das notas de cada uma das questões correspondentes.
A primeira questão discursiva da Formação Geral foi desconsiderada e, portanto, a nota bruta referente às questões discursivas de Formação Geral foi igual à nota da questão mantida, a segunda.
Na componente específica a nota bruta correspondeu à média das notas nas três questões.
As notas brutas das questões objetivas de cada uma das partes, Formação Geral/Componente Específico, foram calculadas a partir dos acertos dos estudantes considerando -se uma regra de três para a conversão para a escala de 0 a 100, descontando -se aquelas questões que foram desconsideradas pela comissão ou que não passaram pelo teste do índice de discriminação ponto-bisserial ( ver seção 1.4.6 para a descrição do índice ).
Questões com índices de discriminação inferior a 0,2 foram excluídas do cômputo da nota do estudante.
As notas brutas de Formação Geral e do Componente Específico para cada um dos estudantes foram calculadas como médias ponderadas das suas notas nas questões objetivas e discursivas, ponderação diferente para cada parte : C\_FG = C\_FG\textasciicircum{}OBJ × 0,6 + C\_FG\textasciicircum{}DIS × 0,4 C\_CE = C\_CE\textasciicircum{}OBJ × 0,85 + C\_CE\textasciicircum{}DIS × 0,15, onde C\_FG é a nota bruta da Formação Geral, C\_FG\textasciicircum{}OBJ é a nota bruta nas questões objetivas da Formação Geral, C\_FG\textasciicircum{}DIS é a nota bruta nas questões discursivas da Formação Geral, C\_CE é a nota bruta do Componente Específico, C\_CE\textasciicircum{}OBJ é a nota bruta nas questões objetivas do Componente Específico, e C\_CE\textasciicircum{}DIS é a nota bruta na parte discursiva do Componente Específico.
A nota bruta da prova, C\_GER, por sua vez, é também uma média ponderada da nota de cada uma das partes, Formação Geral/Componente Específico : C\_GER = 0,25 × C\_FG + 0,75 × C\_CE. 3.1 Estatísticas Básicas da Prova 3.1.1 Estatísticas Básicas Gerais A Tabela 3.1 apresenta as estatísticas básicas da prova por grupo de estudantes.
O total de estudantes inscritos para a prova foi de 3.252.
Destes, 85,9 \% compareceram, sendo que a presença foi maior entre concluintes ( 90,0 \% ) do que entre ingressantes ( 82,6 \% ).
A média geral da prova foi de 43,62 sendo que os ingressantes obtiveram média mais baixa, de 39,72, que os concluintes, cuja média foi 47,98.
O desvio-padrão geral foi de 18,2, sendo o do grupo de ingressantes menor ( 16,3 ), do que o do grupo dos concluintes ( 20,9 ), aqueles possuíram uma distribuição mais homogênea nas notas.
A nota máxima foi 81,0, obtida por pelo menos um concluinte ao passo que a maior nota obtida por um ingressante foi 77,5.
Os coeficientes de assimetria para os dois grupos considerados, o de concluintes e o de ingressantes, bem como para a população como um todo foram todos negativos, indicando uma maior concentração à direita da distribuição e um maior espalhamento à esquerda.
Uma comparação entre o desempenho de ingressantes e concluintes pode ser realizada a partir dos dados apresentados no Gráfico 3.1.
Cumpre destacar que as notas dos concluintes, de um modo geral, foram mais elevadas do que as dos ingressantes, cuja distribuição de notas está à esquerda dos primeiros.
A ideia de melhor desempenho é reforçada pelas notas medianas dos dois grupos : 51,3 entre os concluintes e 45,9 entre os ingressantes.
Nas faixas de notas até 50 pontos o percentual de ingressantes é maior que o dos concluintes, situação esta que se inverte para as faixas de notas superiores a 50 pontos.
A exceção é o primeiro intervalo de 0 a 10 que inclui uma quantidade maior de concluintes do que de ingressantes.
Constatou -se que aproximadamente 20,5 \% dos ingressantes e 46,2 \% dos concluintes situaram -se nas faixas de notas superiores a 50,0 pontos.
O Gráfico 3.2 apresenta informações referentes ao desempenho geral de ingressantes e concluintes, comparando os resultados considerando a Região, Categoria Administrativa e Organização Acadêmica.
Levando -se em conta as notas médias dos estudantes em cada Região, observa -se que há, nos dois grupos, uma variação expressiva entre as Regiões.
As notas médias mais elevadas entre os concluintes ocorreram nas Regiões Sul ( 54,59 ) e Centro-Oeste ( 48,50 ), ambas acima da média nacional ( 47,98 ).
Em relação aos ingressantes, as médias mais altas foram obtidas nas Regiões Sudeste ( 40,74 ) e Nordeste ( 39,97 ) nas demais Regiões a média foi inferior à média nacional ( 39,72 ).
Na Região Centro-Oeste foi obtida a média de notas mais baixa entre os ingressantes ( 35,25 ).
Entre os concluintes, a nota mais baixa foi verificada na Região Norte ( 41,31 ).
Cursos nesta área foram oferecidos em instituições federais, estaduais e privadas.
Levando -se em conta os agrupamentos dos estudantes em categorias administrativas, observa -se que a pontuação média mais elevada foi encontrada em instituições de origem privada ( ingressantes 39,90 ) e estadual ( concluintes 49,57 ) – ambas acima da média nacional.
Tendo como foco as notas médias dos estudantes agrupados por Organização Acadêmica, notamos que as Faculdades apresentaram as maiores médias ( 49,23 para concluintes e 45,15 para ingressantes ), ao passo que as piores notas médias foram observadas nos Centros Universitários ( 42,73 para concluintes e 37,34 para ingressantes ). 3.1.2 Estatísticas Básicas de Formação Geral A Tabela 3.2 apresenta as estatísticas básicas de Formação Geral, por grupo de estudantes.
O número de estudantes inscritos no exame foi de 3.252 e o comparecimento foi de 85,9 \% desse total, sendo que a presença entre concluintes ( 90,0 \% ) foi maior que entre ingressantes ( 82,6 \% ).
A média geral da parte da prova de Formação Geral foi de 42,14, sendo que os ingressantes obtiveram média ligeiramente mais baixa, 41,95, que os concluintes, cuja média foi 42,35.
O desvio-padrão geral foi de 20,5, sendo o do grupo de ingressantes ( 20,1 ), e o do grupo dos concluintes ( 20,9 ), aqueles possuíram uma distribuição mais homogênea nas notas.
A nota máxima foi 100, obtida por estudantes de ambos os grupos.
Para esta parte, assim como para a nota geral, os coeficientes de assimetria são positivos para todos os grupos.
Como, porém, os valores são bem perto de zero, estes indicaram uma ligeira inclinação para a direita.
O Gráfico 3.3 permite a comparação entre concluintes e ingressantes relativamente ao desempenho na parte da prova que aborda a Formação Geral.
Nesta parte, se excluidas as provas com nota zero, tal como na prova como um todo, também os concluintes apresentaram desempenho um pouco superior, embora seja perceptível um maior equilíbrio.
A mediana dos concluintes foi igual à dos ingressantes, enquanto a média dos concluintes foi um pouco maior do que a dos ingressantes.
A razão da discrepância, como já notado para a nota geral esteve relacionada com a maior quantidade de estudantes com nota zero entre os concluintes.
Dos ingressantes, aproximadamente 32,5 \% alcançaram nota superior a 50 ; dentre os concluintes, o percentual dos estudantes nesta faixa de notas foi maior e igual a 37,1 \%.
No Gráfico 3.4 são apresentadas as informações referentes ao desempenho dos ingressantes e dos concluintes na parte de Formação Geral, em diferentes agregações como Região, Categoria Administrativa e Organização Acadêmica.
Considerando -se as notas médias dos estudantes segundo cada Região, observou -se que, tanto no grupo dos ingressantes quanto no grupo dos concluintes, houve diferenças marcantes.
As médias variam, entre os concluintes, de 30,77 na Região Norte a 46,66 na Região Sul.
Entre os ingressantes, os valores foram de 37,12 na Região Norte a 44,98 na Região Sudeste.
Nas Regiões Nordeste e Sul, os concluintes apresentaram médias maiores do que os ingressantes, a situação \textbf{inversa} aconteceu no Norte, Sudeste e Centro-Oeste.
Tendo como foco as categorias administrativas, as maiores notas médias foram das instituições privadas, tanto para concluintes ( 42,87 ) como para ingressantes ( 42,45 ), enquanto às das instituições estaduais foram menores tanto para concluintes ( 40,64 ) e ingressantes ( 38,99 ).
Os estudantes das instituições federais se apresentaram mais homogêneos do que as demais instituições.
Considerando -se o tipo de Organização Acadêmica, Faculdades apresentaram os maiores valores para os concluintes e para os ingressantes, enquanto os Centros Universitários apresentaram os menores valores, também para ambos.
Centros Federais de Educação Tecnológica e Institutos de Educação, Ciência e Tecnologia não estavam representados nesta área.
No Gráfico 3.4 são apresentadas as informações referentes ao desempenho dos ingressantes e dos concluintes na parte de Formação Geral, em diferentes agregações como Região, Categoria Administrativa e Organização Acadêmica.
Considerando -se as notas médias dos estudantes segundo cada Região, observou -se que, tanto no grupo dos ingressantes quanto no grupo dos concluintes, houve diferenças marcantes.
As médias variam, entre os concluintes, de 30,77 na Região Norte a 46,66 na Região Sul.
Entre os ingressantes, os valores foram de 37,12 na Região Norte a 44,98 na Região Sudeste.
Nas Regiões Nordeste e Sul, os concluintes apresentaram médias maiores do que os ingressantes, a situação \textbf{inversa} aconteceu no Norte, Sudeste e Centro-Oeste.
Tendo como foco as categorias administrativas, as maiores notas médias foram das instituições privadas, tanto para concluintes ( 42,87 ) como para ingressantes ( 42,45 ), enquanto às das instituições estaduais foram menores tanto para concluintes ( 40,64 ) e ingressantes ( 38,99 ).
Os estudantes das instituições federais se apresentaram mais homogêneos do que as demais instituições.
Considerando -se o tipo de Organização Acadêmica, Faculdades apresentaram os maiores valores para os concluintes e para os ingressantes, enquanto os Centros Universitários apresentaram os menores valores, também para ambos.
Centros Federais de Educação Tecnológica e Institutos de Educação, Ciência e Tecnologia não estavam representados nesta área. 3.1.3 Estatísticas Básicas em Componente Específico A Tabela 3.3 apresenta as estatísticas básicas referentes ao Componente Específico da prova, tendo sua média geral, de 44,16, menor do que o desempenho da prova de BIBLIOTECONOMIA como um todo, cuja média foi 43,62, como apresentado na Tabela 3.1, que por sua vez foi superior ao desempenho na Formação Geral, que teve média igual a 42,14, conforme mostrado na Tabela 3.2.
Os concluintes apresentaram um desempenho médio ( 49,91 ) superior ao dos ingressantes ( 39,03 ), o mesmo ocorrido para a nota da Formação Geral, havendo ainda indícios de maior homogeneidade entre as notas médias dos ingressantes do que entre as dos concluintes, como sugeriu o desvio padrão, menor no grupo dos ingressantes do que no dos concluintes.
As notas máximas obtidas entre os concluintes e os ingressantes foram iguais a 94,8 e 80,3, respectivamente.
As distribuições de notas para cada grupo se apresentaram com uma inclinação para notas maiores, como evidenciado pelos coeficientes de assimetria negativos.
A situação de assimetria é confirmada pelo Gráfico 3.5 que proporciona uma comparação do desempenho entre concluintes e ingressantes em relação ao Componente Específico da prova.
Claramente a distribuição das notas médias dos ingressantes, à direita daquela dos concluintes, evidenciou o percentual mais alto de notas médias nas faixas mais baixas para este grupo, com exceção do primeiro, que como já comentado, incluiu uma maior proporção de zeros entre os concluintes.
As medianas das distribuições das notas dos ingressantes e dos concluintes foram 39 e 51, respectivamente, enfatizando o melhor desempenho dos concluintes nesta parte da prova.
Entre os concluintes, 58,3 \% e entre os ingressantes, 37,7 \%, obtiveram notas acima de 50.
Assim como os Gráficos 3.2 e 3.4, o Gráfico 3.6 apresenta uma comparação dos resultados em relação às Regiões, Categoria Administrativa e Organização Acadêmica, agora levando em conta o desempenho dos ingressantes e dos concluintes no Componente Específico da prova.
Tendo como foco as notas médias dos estudantes, em determinada Região, observa -se que, entre os concluintes, as mais elevadas foram encontradas nas Regiões Sul ( 57,29 ), Nordeste ( 50,81 ) e Centro-Oeste ( 50,27 ) valores acima da média nacional para esta parte ( 49,90 ).
Entre os ingressantes, as maiores notas ocorreram nas Regiões : Nordeste ( 40,41 ), Sudeste ( 39,37 ) e Sul ( 39,36 ), todos acima da média nacional ( 39,03 ).
Novamente, para a análise segundo a Categoria Administrativa, somente entidades municipais não foram representadas.
As médias mais elevadas entre os concluintes foram encontradas em instituições de origem estadual ( 52,59 ) e entre os ingressantes, foram encontradas em instituições de origem federal e privada ( 39,10 para ambas ).
Quanto à Organização Acadêmica, o comentário é o mesmo já feito para a nota bruta e referente à Formação Geral : universidades, centros universitários e faculdades foram representadas neste curso.
Considerando -se o tipo de Organização Acadêmica, faculdades apresentaram os maiores valores, seguidas das universidades e dos centros universitários.
Assim como os Gráficos 3.2 e 3.4, o Gráfico 3.6 apresenta uma comparação dos resultados em relação às Regiões, Categoria Administrativa e Organização Acadêmica, agora levando em conta o desempenho dos ingressantes e dos concluintes no Componente Específico da prova.
Tendo como foco as notas médias dos estudantes, em determinada Região, observa -se que, entre os concluintes, as mais elevadas foram encontradas nas Regiões Sul ( 57,29 ), Nordeste ( 50,81 ) e Centro-Oeste ( 50,27 ) valores acima da média nacional para esta parte ( 49,90 ).
Entre os ingressantes, as maiores notas ocorreram nas Regiões : Nordeste ( 40,41 ), Sudeste ( 39,37 ) e Sul ( 39,36 ), todos acima da média nacional ( 39,03 ).
Novamente, para a análise segundo a Categoria Administrativa, somente entidades municipais não foram representadas.
As médias mais elevadas entre os concluintes foram encontradas em instituições de origem estadual ( 52,59 ) e entre os ingressantes, foram encontradas em instituições de origem federal e privada ( 39,10 para ambas ).
Quanto à Organização Acadêmica, o comentário é o mesmo já feito para a nota bruta e referente à Formação Geral : universidades, centros universitários e faculdades foram representadas neste curso.
Considerando -se o tipo de Organização Acadêmica, faculdades apresentaram os maiores valores, seguidas das universidades e dos centros universitários. 3.2 Análise das Questões Objetivas 3.2.1 Formação Geral A Tabela 3.4 apresenta as estatísticas básicas relativas às oito questões objetivas da Formação Geral dos estudantes.
Adiantamos que a primeira questão foi desconsiderada pelo critério da correlação ponto-bisserial.
Como pode ser observado na Tabela 3.4, os estudantes concluintes obtiveram desempenho ligeiramente inferior, com nota média de 55,89, em relação aos estudantes ingressantes cuja nota média foi 56,31.
A variabilidade entre os concluintes foi cerca de 0,5 maior do que entre os ingressantes, como indicaram os desvios-padrão de 23,2, para o grupo de ingressantes e 22,7 para o grupo dos concluintes.
A amplitude total das notas foi a mesma para ambos os grupos : nota mínima de 0 e máxima de 100.
A mediana foi igual para ambos os grupos : 57,1.
A assimetria, medida pelo coeficiente de assimetria ( skewness ) foi um pouco menor em módulo para os ingressantes ( -0,3 ) do que para os concluintes ( -0,4 ), ainda que ambos fossem negativos indicando uma assimetria para a esquerda, fato reforçado pelos valores menores das médias dos dois grupos quando comparados à mediana correspondente.
A Tabela 3.5 apresenta a classificação das questões objetivas da Formação Geral, segundo o índice de facilidade ( questões em vermelho foram desconsideradas pelo critério ponto-bisserial ).
Como as questões foram comuns a todas as áreas, a tabela apresenta a classificação considerando -se os concluintes de todas as áreas e aqueles do curso de BIBLIOTECONOMIA.
De acordo com os índices obtidos, as questões objetivas da prova foram assim avaliadas para a área de BIBLIOTECONOMIA : das oito questões, três tiveram A Tabela 3.6 descreve os resultados das análises das questões objetivas relativas à Formação Geral, segundo o poder de discriminação, utilizando -se para tal o índice de discriminação ponto-bisserial.
Como as questões foram comuns a todas as áreas, a tabela apresenta a classificação considerando -se os concluintes de todas as áreas e aqueles do curso de BIBLIOTECONOMIA.
A desconsideração da questão segundo este critério considerou o universo de todos os concluintes.
Uma das oito questões apresentou índice igual ou acima de 0,40 ( nenhuma para a população de concluintes como um todo ), sendo, portanto, classificada como muito boa nesta característica, para este grupo de estudantes, duas foram classificadas como boas com índice entre 0,30, inclusive, e 0,40 ( três para os concluintes como um todo ) e quatro foram classificadas como médias, com índice entre 0,20, inclusive, e 0,30 ( quatro para os concluintes como um todo ).
Uma das questões foi classificada como fraca tanto para a população de concluintes como um todo como para a sub-população de concluintes de BIBLIOTECONOMIA e foi desconsiderada pelo critério. entre 60 \%, inclusive, e 85 \% de acertos e foram classificadas como fáceis ( somente uma para o total de concluintes ) ; outras duas foram consideradas medianas, alcançando entre 40 \%, inclusive, e 60 \% de acertos ( três para o total de concluintes ) ; duas obtiveram entre 15 \%, inclusive, e 40 \% de acertos, sendo consideradas difíceis ( três para o total de concluintes ).
Por fim, uma foi classificada como muito difícil tento obtido percentual menor do que 15 \% de acertos ( uma para o total de concluintes ), justamente a questão desconsiderada pelo critério ponto-bisserial.
Nenhuma questão foi classificada na categoria muito fácil.
Em linhas gerais podemos dizer que os concluintes de BIBLIOTECONOMIA tiveram um melhor desempenho nesta parte do que o total da população de concluintes de todos os cursos avaliados em 2009.
A Tabela 3.6 descreve os resultados das análises das questões objetivas relativas à Formação Geral, segundo o poder de discriminação, utilizando -se para tal o índice de discriminação ponto-bisserial.
Como as questões foram comuns a todas as áreas, a tabela apresenta a classificação considerando -se os concluintes de todas as áreas e aqueles do curso de BIBLIOTECONOMIA.
A desconsideração da questão segundo este critério considerou o universo de todos os concluintes.
Uma das oito questões apresentou índice igual ou acima de 0,40 ( nenhuma para a população de concluintes como um todo ), sendo, portanto, classificada como muito boa nesta característica, para este grupo de estudantes, duas foram classificadas como boas com índice entre 0,30, inclusive, e 0,40 ( três para os concluintes como um todo ) e quatro foram classificadas como médias, com índice entre 0,20, inclusive, e 0,30 ( quatro para os concluintes como um todo ).
Uma das questões foi classificada como fraca tanto para a população de concluintes como um todo como para a sub-população de concluintes de BIBLIOTECONOMIA e foi desconsiderada pelo critério.
Na sequência, no Quadro 3.1, estão especificadas, por questão, as habilidades aferidas pelas questões de Formação Geral da prova da BIBLIOTECONOMIA, bem como os resultados da análise relativa aos índices de facilidade e de discriminação de cada questão.
O índice de facilidade para a área de BIBLIOTECONOMIA teve uma variação considerável, oscilando entre 0,11 até 0,85 ( na população de concluintes o intervalo foi de 0,11 a 0,85 ) e o de discriminação ficou entre 0,09 e 0,46.
A questão 8, que aferiu as habilidades de “ Ler e interpretar textos ”, “ Analisar e criticar informações ” ; ” Extrair conclusões por indução e/ou dedução ”, ” Estabelecer relações, comparações e contrastes em diferentes situações ”, “ Detectar contradições ”, “ Fazer escolhas valorativas avaliando consequências ”, “ Questionar a realidade ” e “ Argumentar coerentemente ” foi a que apresentou, para a população de concluintes de BIBLIOTECONOMIA, maior poder discriminatório com índice de 0,46.
A questão mais fácil, com 85 \% de acertos, foi também a de número 8.
De outra parte, a questão com menor índice de facilidade e de discriminação em Formação Geral foi a questão 1, que foi desconsiderada, e que deveria aferir as habilidades de “ Ler e interpretar textos ”, “ Analisar e criticar informações ” ; ” Extrair conclusões por indução e/ou dedução ”, ” Estabelecer relações, comparações e contrastes em diferentes situações ”, “ Detectar contradições ”, “ Fazer escolhas valorativas avaliando consequências ”, “ Questionar a realidade ” e “ Argumentar coerentemente ”.
O Gráfico 3.7 mostra, a título de ilustração, o comportamento das respostas à questão 3 de Formação Geral para os estudantes concluintes da área de BIBLIOTECONOMIA como função da nota geral na prova.
Trata -se de uma questão considerada de dificuldade média para os concluintes de BIBLIOTECONOMIA e também para o conjunto de concluintes de todas as áreas.
O poder de discriminação foi médio para os estudantes da área e para o conjunto de todos os concluintes.
Neste Gráfico, cada uma das cinco curvas, representou o percentual de respostas em determinada alternativa da questão, em função do escore total dos estudantes concluintes da área de BIBLIOTECONOMIA.
A curva roxa corresponde à alternativa E, a correta para esta questão.
Assim, observa -se que entre os estudantes com notas mais baixas – desempenho geral mais fraco na prova como um todo –, a tendência foi a escolha de alternativas incorretas, especialmente a alternativa C – curva vermelha, seguida da alternativa D – curva em verde escuro.
A proporção de estudantes que selecionou a resposta correta aumentou com o desempenho deste grupo de estudantes chegando ao acerto de 100 \% dos estudantes com notas pelo menos 87,5.
Esta análise permitiu verificar como a questão discriminou entre os grupos de desempenho de acordo com o índice obtido.
Os Gráficos relativos às demais questões de Formação Geral constam do Anexo I. 3.2.2 Componente Específico A Tabela 3.7 apresenta as estatísticas básicas em relação às questões objetivas do Componente Específico.
Como pode ser observado, os estudantes concluintes obtiveram desempenho superior, com uma média de 52,48 pontos, em relação aos estudantes ingressantes com uma média de 41,17 pontos.
Os desvios-padrão indicaram que a variabilidade entre os concluintes foi superior ( 18,1 ) do que entre os ingressantes ( 14,9 ).
Os coeficientes de assimetria para os dois grupos foram negativos indicando uma inclinação para a direita.
A assimetria foi, porém, maior em módulo entre os concluintes.
A análise das questões objetivas do Componente Específico em BLIOTECONOMIA, em relação ao grau de facilidade da prova ( ver Tabela 3.8 ), mostrou que, das 27 questões, 14 encontraram -se na categoria difícil ( entre 15 \% e 40 \%, não incluso, de acertos ), 4 são de nível médio ( entre 40 \% e 60 \%, não incluso, de acertos ), 6 estão na categoria fácil ( entre 60 \% e 85 \%, não incluso, de acertos ) e nenhuma teve a classificação muito fácil ( a partir de 85 \% de acertos ).
Nenhuma questão teve percentual de acertos menor do que 15 \% ( classificação muito difícil ) e 3 questões foram desconsideradas pela Comissão Assessora da Área de BIBLIOTECONOMIA.
Das 27 questões objetivas relativas ao Componente Específico ( ver Tabela 3.9 ), 4 obtiveram o índice de discriminação muito bom, 9 obtiveram o índice de discriminação bom e 6 tiveram nível de discriminação médio.
O índice fraco que eliminaria a questão, foi alcançado por 5 questões.
Foram desconsideradas pela comissão 3 questões.
Na sequência, no Quadro 3.2 estão especificadas, por questão, as habilidades aferidas pelas questões do Componente Específico da prova, bem como os resultados da análise relativa aos índices de facilidade e de discriminação de cada questão.
O Gráfico 3.8 mostra, a título de ilustração, o comportamento da questão 11 do Componente Específico para os estudantes concluintes da área de BIBLIOTECONOMIA como função da nota geral na prova.
Trata -se de uma questão considerada difícil para os concluintes de BIBLIOTECONOMIA.
O poder de discriminação foi médio para os estudantes da área.
Neste Gráfico, cada uma das cinco curvas, representa o percentual de respostas em determinada alternativa da questão, em função da nota total dos estudantes concluintes de BIBLIOTECONOMIA.
A curva verde escuro correspondeu à alternativa D, a correta para esta questão.
Assim, observou -se que entre os estudantes com escores mais baixos – desempenho geral mais fraco na prova como um todo –, a tendência foi a escolha de uma alternativa incorreta, a alternativa B – curva verde claro ou a alternativa C – curva vermelha.
A proporção de estudantes que selecionou a resposta correta aumentou com o desempenho deste grupo de estudantes chegando ao acerto de 100 \% dos estudantes com notas pelo menos 87,5.
Esta análise permitiu verificar como a questão discriminou entre os grupos de desempenho.
Os Gráficos relativos às demais questões do Componente Específico constam do Anexo I. 3.3 Análise das Questões Discursivas 3.3.1 Formação Geral A Tabela 3.10 apresenta as estatísticas básicas da prova por grupo de estudantes.
A análise dos resultados de desempenho dos estudantes nas questões discursivas que são apresentados na tabela 3.10 e no gráfico 3.9, mostrou que as notas foram mais baixas no conjunto dessas questões que no das questões objetivas.
O total de estudantes inscritos para a prova foi de 3.252.
Destes, 85,9 \% compareceram, sendo que a presença foi maior entre concluintes ( 90,0 \% ) do que entre ingressantes ( 82,6 \% ).
A média geral da prova foi de 21,28, sendo que os ingressantes obtiveram média mais baixa, de 20,51, que os concluintes, cuja média foi 22,15.
O desvio-padrão geral foi de 33,7, sendo o do grupo de ingressantes ligeiramente menor ( 33,5 ), do que o do grupo dos concluintes ( 33,8 ), aqueles possuem uma distribuição mais homogênea nas notas.
A nota máxima foi 100,0, atingida por concluintes e ingressantes.
Os coeficientes de assimetria para os dois grupos considerados, o de concluintes e o de ingressantes, bem como para a população como um todo foram todos positivos, indicando uma maior concentração à esquerda da distribuição.
Para todos os grupos mais da metade dos estudantes presentes obtiveram nota zero nas questões discursivas da Formação Geral.
Das 1.891 provas com nota zero, 38 ( 2 \% ) continham um protesto, 329 ( 17,4 \% ) estavam em branco e 1524 ( 80,6 \% ) foram respondidas, porém de forma inadequada.
Entre os concluintes cerca de 81 \% dentre aqueles com nota zero estavam com a prova respondida de forma incorreta.
Entre os ingressantes 80 \% das notas nulas ( zero ) decorreram do desempenho do próprio estudante verificadas na correção.
Uma comparação entre o desempenho de ingressantes e concluintes pode ser realizada a partir dos dados apresentados no Gráfico 3.9.
As notas dos concluintes, de um modo geral, foram ligeiramente mais elevadas do que as dos ingressantes, cuja distribuição de notas está à direita dos primeiros.
A ideia de melhor desempenho foi reforçada pela proporção de notas nulas : 66,4 \% entre os concluintes e 68,8 \% entre os ingressantes.
Entre os concluintes, 25,0 \% obtiveram notas superiores a 50 e entre os ingressantes, 22,8 \%.
Na sequência, discorrer -se-á, especificamente, sobre uma das questões discursivas de Formação Geral, a número 10 já que a número 9 foi desconsiderada.
Em especial, serão analisados os conteúdos englobados na questão e os resultados obtidos pelos estudantes.
Ao final, são apresentados os comentários da banca corretora a respeito do observado nas respostas dos estudantes, suas impressões e conclusões.
Cumpre esclarecer que, tendo em vista que as questões discursivas de Formação Geral foram padronizadas, ou seja, constam de todas as provas, os comentários da banca foram os mesmos para todas as áreas acadêmicas, eis que direcionados a todos os estudantes que participaram do ENADE/2009. 3.3.1.1 Análise da Questão Discursiva 9 de Formação Geral DESCONSIDERADA 3.3.1.2 Análise da Questão Discursiva 10 de Formação Geral A Questão Discursiva nº. 10 buscou avaliar as habilidades : • Ler e interpretar textos ; • Analisar e criticar informações ; • Extrair conclusões por indução e/ou dedução ; • Estabelecer relações, comparações e contrastes em diferentes situações ; • Fazer escolhas valorativas avaliando consequências ; • Questionar a realidade ; • Argumentar coerentemente.
O Quadro 3.3 apresenta as habilidades/competências, temas predominantes e estatísticas básicas da questão 10 de Formação Geral por grupo de estudantes.
Cumpre notar que, como a questão 9 foi desconsiderada, a nota referente às questões discursivas da Formação Geral é idêntica à nota da questão 10.
A média geral desta questão foi de 21,28, sendo que os ingressantes obtiveram média mais baixa, de 20,51, que os concluintes, cuja média foi 22,15.
O desvio-padrão geral foi de 33,7, sendo que o do grupo de ingressantes foi ligeiramente menor ( 33,5 ), em relação ao grupo dos concluintes ( 33,8 ), aqueles possuíram uma distribuição mais homogênea nas notas.
A nota máxima foi 100,0, atingida por concluintes e ingressantes.
Os coeficientes de assimetria para os dois grupos considerados, o de concluintes e o de ingressantes, bem como para a população como um todo foram positivos, indicando uma maior concentração à esquerda da distribuição.
Para todos os grupos, mais da metade dos estudantes presentes tiveram nota zero nesta questão discursiva de Formação Geral.
Das 1.891 provas com nota zero, 38 ( 2,0 \% ) continham um protesto, 329 ( 17,4 \% ) estavam em branco e 1.524 ( 80,6 \% ) foram respondidas, porém de forma inadequada.
Entre os concluintes a maioria ( 80,9 \% ) dentre aqueles com nota zero estava com a prova respondida incorretamente.
Entre os ingressantes 80 \% das notas nulas ( zero ) decorreram do desempenho do próprio estudante verificadas na correção.
Uma comparação entre o desempenho de ingressantes e concluintes para esta questão foi exatamente igual à já realizada para a nota discursiva em Formação Geral a partir dos dados apresentados no Gráfico 3.9. 3.3.1.3 Comentários sobre as respostas à Questão Discursiva 10 Apresentação A questão 10 de Formação Geral apresentou um questionamento capaz de analisar a habilidade do estudante em reconhecer o que é um direito social, optar por um direito social vigente que possa diminuir as desigualdades de renda familiar no Brasil e discutir os argumentos em favor do direito social escolhido como suporte à resposta.
Para análise das respostas apresentadas pelos estudantes, houve uma leitura inicial bastante criteriosa.
Foram considerados para efeito de pontuação, os seguintes argumentos : • Acesso à educação pública, gratuita e de qualidade, garantindo ao cidadão melhores condições de vida ; • Assistência à saúde, em seu contexto mais amplo, garantindo os cuidados essenciais de saúde, baseados em técnicas apropriadas ao alcance de todas as pessoas, contando com a participação da população ; • Permanência do estudante na escola, em todos os níveis escolares ( da educação infantil a educação superior, possibilitando a qualificação profissional ), o que garante melhores condições de trabalho e remuneração.
Considerou -se a qualificação e o desempenho docente, desenvolvendo nos estudantes melhores expectativas de aprendizagem e interesse pelo conteúdo da escola.
Os professores desenvolvendo metodologias adequadas e conhecimento pleno quanto a questões de aprendizagem ; • Condições dignas de trabalho, com remuneração que garanta qualidade de vida.
Foram consideradas as condições de trabalho, a preparação profissional para o trabalho e os cursos de aperfeiçoamento, oportunizando momentos de aprendizagem e desenvolvimento, a partir de situações desencadeadoras ; • Ser proprietário do imóvel que reside, reduzindo os gastos com aluguel, garantindo melhores condições de vida, equilíbrio financeiro e construção do bem-estar familiar.
Desempenho dos estudantes na questão É importante salientar que havia respostas que desenvolveram adequadamente o questionamento apresentado, com fundamentação baseada em argumentos coerentes.
Contudo, em um número considerável de respostas, havia apenas um argumento e não dois, conforme fora solicitado.
Muitos estudantes informam sobre problemas da educação, de um modo geral.
Algumas observações mencionadas pelos corretores durante o processo de correção foram : 3.3.2 Componente Específico Nas questões discursivas do Componente Específico – questões de números 38, 39 e 40 –, a diferença de desempenho entre os grupos não foi muito acentuada, inclusive menor do que para as questões discursivas de Formação Geral.
A Tabela 3.11 contém as estatísticas da média destas questões.
Observa -se que no grupo dos concluintes a nota média alcançada foi 35,62, enquanto que no grupo dos ingressantes foi de 27,18, ambas maiores do que as obtidas nas questões discursivas de Formação Geral, que foram de 20,51 para os ingressantes e 22,15 para os concluintes ( Tabela 3.10 ).
Contudo, as médias foram menores dos que as obtidas para as questões objetivas do Componente Específico ( Tabela 3.7 ) : respectivamente 41,17 e 52,48 para ingressantes e concluintes.
A análise de cada uma destas questões será feita nas subseções seguintes. • há preocupação com uma educação de qualidade, com sugestões de cursos profissionalizantes e métodos preparatórios de aulas. • a educação é revelada como de baixa qualidade, exaltando -se treinamentos e cursos de atualização para a classe docente. • a saúde como essência para a vida.
O SUS é colocado como parte integrante de direito social, garantindo às pessoas fácil acesso. • oportunidade de emprego como forma de ascensão social, direito, cidadania. • os políticos brasileiros como causadores do despreparo educacional.
As respostas analisadas apresentaram norma culta, coesão e coerência razoáveis, que garantiram a compreensão da mensagem desenvolvida pelos estudantes.
Os erros ortográficos foram poucos e não comprometeram a discussão sobre o assunto e sobre os pontos de vista abordados.
Quanto à coerência, houve um desvio em relação ao objetivo da questão, já que alguns estudantes, através dos seus argumentos retrataram reivindicações, reclamações e situações de problemas enfrentados pelos brasileiros, citando golpes políticos.
Houve erros quanto à concordância verbal e nominal.
Houve, também, dificuldade em utilizar os conectores corretamente.
Alguns estudantes apresentaram em suas respostas redundância e prolixidade.
O desempenho foi considerado satisfatório, pois a maioria dos estudantes soube responder ao questionamento.
Houve domínio do conteúdo, pois as ideias, os argumentos e a originalidade corresponderam ao esperado. 3.3.2 Componente Específico Nas questões discursivas do Componente Específico – questões de números 38, 39 e 40 –, a diferença de desempenho entre os grupos não foi muito acentuada, inclusive menor do que para as questões discursivas de Formação Geral.
A Tabela 3.11 contém as estatísticas da média destas questões.
Observa -se que no grupo dos concluintes a nota média alcançada foi 35,62, enquanto que no grupo dos ingressantes foi de 27,18, ambas maiores do que as obtidas nas questões discursivas de Formação Geral, que foram de 20,51 para os ingressantes e 22,15 para os concluintes ( Tabela 3.10 ).
Contudo, as médias foram menores dos que as obtidas para as questões objetivas do Componente Específico ( Tabela 3.7 ) : respectivamente 41,17 e 52,48 para ingressantes e concluintes.
A análise de cada uma destas questões será feita nas subseções seguintes.
Os valores positivos para a medida de assimetria indicaram uma distribuição com cauda para a direita.
Dos 2.794 presentes, 30,5 \% obtiveram nota zero.
A tabela 3.12 apresenta um detalhamento da situação destes estudantes.
Cerca de 42 \% não fez nenhuma das 3 questões.
Esta proporção foi ligeiramente maior entre os concluintes ( 43 \% ) do que entre os ingressantes ( 41,3 \% ).
A combinação de questões em branco e preenchidas correspondeu a 93,2 \% das provas.
Menos de 5 \% no total das provas correspondeu a 3 questões com protesto.
O Gráfico 3.10 apresenta a distribuição das notas nas questões discursivas do Componente Específico por grupo de estudantes.
A moda se situou no primeiro intervalo, aquele com valores entre 0 e 10 incluindo os extremos do intervalo.
A impressão de assimetria informada pelos coeficientes de assimetria positivos ( Tabela 3.11 ) foi reforçada no gráfico pela concentração no primeiro intervalo. 3.3.2.1 Análise da Questão Discursiva 38 do Componente Específico A questão 38, cujas habilidades aferidas encontram -se descritas no Quadro 3.4, que ainda contém outras informações, seguido da Tabela 3.13 e do Gráfico 3.11, foi aquela em que os grupos de estudantes observados obtiveram um desempenho intermediário entre as questões discursivas do Componente Específico.
A nota média do grupo dos ingressantes foi de 28,56, e a dos concluintes 36,74.
O percentual de notas zero obtidas nos dois grupos foi de 39,5 \%, maior entre os ingressantes ( 43,5 \% ) do que entre os concluintes ( 34.9 \% ).
Cerca de 40 \% dos concluintes e dos ingressantes que tiraram zero deixaram a questão em branco ( ver Tabela 3.13 ).
Uma pequena minoria ( 5,6 \% ) utilizou a questão para protestar.
Apesar do baixo desempenho, em ambos os grupos a nota máxima de 100,0 pontos foi alcançada.
O Gráfico 3.11 apresenta a distribuição das notas na questão 38 discursiva do Componente Específico por grupo de estudantes.
A moda se situou no primeiro intervalo, aquele com valores entre 0 e 10 incluindo os extremos do intervalo.
A impressão de assimetria informada pelos coeficientes de assimetria positivos ( Quadro 3.4 ) foi reforçada no gráfico pela concentração no primeiro intervalo.
A distribuição apresentou outros máximos locais no sétimo e no décimo intervalos para concluintes e no terceiro e no quarto para ingressantes. 3.3.2.2 Comentários sobre as respostas à Questão Discursiva 38 A questão teve como objetivo envolver conteúdos relativos aos fundamentos teóricos da biblioteconomia e da ciência da informação, tendo como referência o perfil explicitado nas diretrizes curriculares da área, especialmente no aspecto referente à reflexão crítica sobre seu papel social.
A questão abordou ainda o desempenho de atividades de identificação das necessidades informacionais de indivíduos, de grupos e de comunidades.
A correção das ações apontadas foi feita, sendo que para cada dificuldade, havia duas ações solicitadas, considerando as opções do padrão de resposta elaborado pela equipe.
O nível foi considerado baixo.
Considerando o desempenho dos estudantes, enfatizamos que aproximadamente mais da metade ficou por aproximação na faixa de 21 a 60 pontos.
No que se refere ao domínio do conteúdo, a questão visava aferir o conhecimento quanto aos fundamentos teóricos da biblioteconomia e da ciência da informação, o que os estudantes resguardaram em suas respostas, de forma generalizada.
No entanto, a maioria dos estudantes enfatizou o estudo de usuários e de comunidades como os serviços básicos a serem oferecidos pela biblioteca.
Poucos abordaram aspectos relacionados aos serviços técnicos.
Como a questão apresentou duas ações em sua preposição, conforme padrão de respostas, a correção observou tal conjuntura para pontuar a expressividade dos estudantes.
É relevante salientar que um pequeno grupo de estudantes que apresentou conteúdo claro e coerente, respeitou a norma culta. 3.3.2.3 Análise da Questão Discursiva 39 do Componente Específico A questão 39, cujas habilidades aferidas encontram -se descritas no Quadro 3.5, que ainda contém outras informações, seguido da Tabela 3.14 e do Gráfico 3.12, foi aquela em que os grupos de estudantes observados obtiveram o pior desempenho entre as questões discursivas do Componente Específico.
A nota média do grupo dos ingressantes foi de 13,95, e a dos concluintes 25,94, bem maior entre os concluintes.
Isto ocorreu, principalmente pelo grande percentual de notas nulas : 61,4 \% entre os ingressantes e 48,1 \% entre os concluintes.
Das provas com notas nulas nesta questão, uma grande proporção foi deixada em branco : 45,4 \% entre os ingressantes e 43,2 \% entre os concluintes ( ver Tabela 3.14 ).
Entre os com nota nula, uma minoria ( 8 \% ) utilizou a questão para protestar.
Apesar do baixo desempenho, em ambos os grupos a nota máxima de 100,0 pontos foi alcançada.
Os coeficientes de assimetria positivos indicaram uma concentração à direita e espalhamento à esquerda.
O Gráfico 3.12 apresenta a distribuição das notas na questão 39 discursiva do Componente Específico por grupo de estudantes.
A moda se situou no primeiro intervalo, aquele com valores entre 0 e 10 incluindo os extremos do intervalo.
A impressão de assimetria informada pelos coeficientes de assimetria positivos ( Quadro 3.5 ) foi reforçada no gráfico pela concentração no primeiro intervalo.
A distribuição apresentou outros máximos locais : no segundo e último intervalos para concluintes e no segundo e terceiro intervalos para ingressantes. 3.3.2.4 Comentários sobre as respostas à Questão Discursiva 39 A questão apresentou como objetivo avaliar o conhecimento sobre recursos e serviços de informação oferecidos em uma unidade de serviços informacionais, através da explicitação de um contexto de uso de informações especializadas.
Com grau de dificuldade alto, esta foi a questão em que a abordagem dos estudantes foi mais cautelosa.
Quanto ao desempenho dos estudantes, observou -se que esta questão foi a que teve menor média das notas das respostas, tanto em relação aos estudantes concluintes, quanto em relação aos estudantes ingressantes.
Foi também a questão que apresentou o maior número de respostas em branco.
Pode -se observar que houve uma tendência dos estudantes a uma divagação em torno de temas mais genéricos, de se valorizar as necessidades dos usuários, sem se aprofundar muito na resposta solicitada.
A maioria dos estudantes enfatizou o estudo de usuários e de comunidades como os serviços básicos a serem oferecidos pela biblioteca.
Poucos abordaram aspectos relacionados aos serviços técnicos.
Observamos que alguns estudantes apresentaram texto com conteúdo adequado, claro e coerente, o que demonstrou o entendimento da questão em relação ao enunciado. 3.3.2.5 Análise da Questão Discursiva 40 do Componente Específico A questão 40, cujas habilidades aferidas encontram -se descritas no Quadro 3.6, que ainda contém outras informações, seguido da Tabela 3.15 e do Gráfico 3.13.
A questão foi aquela em que os grupos de estudantes observados obtiveram o melhor desempenho entre as questões discursivas do Componente Específico.
A nota média do grupo dos ingressantes foi de 39,11 e a dos concluintes 44,25.
Nesta questão ocorreu também o menor percentual de notas nulas : 40,6 \% entre os ingressantes e 34,4 \% entre os concluintes.
Aqui também a maioria das provas com nota zero decorreu de provas deixadas em branco : 47,1 \% entre os ingressantes e 45,8 \% entre os concluintes ( ver Tabela 3.15 ).
Uma pequena minoria ( 6 \% ), dentre aqueles com nota zero, utilizou a questão para protestar.
Em ambos os grupos a nota máxima de 100,0 pontos foi alcançada.
Os coeficientes de assimetria positivos indicaram uma inclinação para a esquerda.
O Gráfico 3.13 apresenta a distribuição das notas na questão 40 discursiva do Componente Específico por grupo de estudantes.
A moda se situou no primeiro intervalo, aquele com valores entre 0 e 10 incluindo os extremos do intervalo.
A impressão de assimetria informada pelos coeficientes de assimetria positivos ( Quadro 3.6 ) foi reforçada no gráfico pela concentração no primeiro intervalo.
A distribuição apresentou outros máximos locais : no sétimo, nono e último intervalos para ingressantes e concluintes. 3.3.2.6 Comentários sobre as respostas à Questão Discursiva 40 Foi solicitado aos estudantes uma análise de forma coerente e contextualizada da situação colocada, demonstrando, assim, seus conhecimentos sobre gestão e em recursos informacionais oferecidos a uma biblioteca voltada para o atendimento do público em geral.
O nível de dificuldade foi mediano, para que o alcance dos estudantes fosse relevante.
Esta foi a questão que teve o maior número de respostas dos estudantes que mostraram dominar e valorizar a questão social e cultural da informação em comunidades que apresentam características de supressão de aparatos e unidades informacionais.
A incidência em concordar com a questão pela escolha do serviço e as razões apresentadas foram bem fundamentadas.
A maioria dos estudantes enfatizou o estudo de usuários e de comunidades como os serviços básicos a serem oferecidos pela biblioteca.
Apesar de não ser a ênfase maior da questão, poucos abordaram aspectos relacionados aos serviços técnicos.

% matched lemmas: inverso
\end{textsample}
